% 核心结论：修复分流电阻被导线短接以及表头到分压电阻的断线
% 草图：[真实几何与节点 | 条件、推导、结论]；17cm印宽，最小字号8pt。
% 图契约：表头改装：并联分流与串联分压；图中的端口、几何、箭头和方程对应本题。
% 版面：沿用前三模块的几何图+推导分区；目标印宽17cm，简图至中等密度。
% 布局：左侧真实结构，右侧关键关系；真实电路边不作为架构流程箭头处理。
\documentclass[tikz,border=7pt]{standalone}
\usepackage{fontspec,xeCJK,amsmath,amssymb}
\IfFontExistsTF{Times New Roman}{\setmainfont{Times New Roman}}{\setmainfont{TeX Gyre Termes}}
\IfFontExistsTF{Heiti SC}{\setCJKmainfont[AutoFakeBold=1.5]{Heiti SC}}{\setCJKmainfont[AutoFakeBold=1.5]{FandolHei-Regular}}
\usetikzlibrary{arrows.meta,bending,calc,patterns,positioning,decorations.pathmorphing}
\definecolor{ink}{HTML}{243745}
\definecolor{blue}{HTML}{275DAB}
\definecolor{green}{HTML}{177C63}
\definecolor{red}{HTML}{BD4935}
\definecolor{orange}{HTML}{AA6C16}
\definecolor{purple}{HTML}{715393}
\tikzset{every node/.style={font=\fontsize{9.3}{12.5}\selectfont,text=ink},
arrow/.style={-{Stealth[length=5pt 1.5,width'=0pt 3]},shorten >=2pt,shorten <=1pt,line width=1pt},
arrow short/.style={-{Stealth[length=3pt 1.0,width'=0pt 2.5]},shorten >=1pt,shorten <=0pt,line width=.8pt},
axis/.style={arrow,draw=ink!60,line width=.6pt},
vec/.style={arrow,draw=blue},
wire/.style={draw=ink!75,line width=.8pt},
guide/.style={dashed,draw=ink!45,line width=.5pt},
card/.style={rounded corners=2pt,draw=blue!65,fill=blue!4,align=center,text width=6.5cm,inner xsep=9pt,inner ysep=7pt},
title/.style={font=\fontsize{10.5}{14}\selectfont\bfseries,text width=15cm,align=center},
label/.style={font=\fontsize{8.4}{11}\selectfont,align=center},
note/.style={label,text width=6.1cm}}
\begin{document}
\begin{tikzpicture}[x=1cm,y=1cm]
\path[use as bounding box] (-8,-3.5) rectangle(8,3.65);
\node[title] at (0,3.25) {\quad 表头改装：并联分流与串联分压};
\begin{scope}[shift={(-4.05,0)}]
\node[label]at(0,2.4){电流表：表头与分流电阻并联};\draw[wire,fill=white](0,1.3)circle(0.3);\node at(0,1.3){$G$};
\draw[wire](-2.5,1.3)--(-.3,1.3)(.3,1.3)--(2.5,1.3)(-2,1.3)--(-2,.1)(2,1.3)--(2,.1);\draw[wire,draw=ink](-2,0.1)--(-0.275,0.1) (0.275,0.1)--(2,0.1);\begin{scope}[shift={(0.0,0.1)},rotate=0.0]\draw[draw=ink,fill=white,line width=.8pt](-0.275,-.085)rectangle(0.275,.085);\end{scope}
\node[label]at(0.0,0.52){$R_s$};
\fill[ink] (-2,1.3)circle(2pt);
\fill[ink] (2,1.3)circle(2pt);
\node[label]at(0,-.75){电压表：表头与分压电阻串联};\draw[wire,fill=white](-1.25,-1.65)circle(0.3);\node at(-1.25,-1.65){$G$};
\draw[wire](-2.5,-1.65)--(-1.55,-1.65)(-.95,-1.65)--(-.5,-1.65);\draw[wire,draw=ink](-0.5,-1.65)--(0.725,-1.65) (1.275,-1.65)--(2.5,-1.65);\begin{scope}[shift={(1.0,-1.65)},rotate=0.0]\draw[draw=ink,fill=white,line width=.8pt](-0.275,-.085)rectangle(0.275,.085);\end{scope}
\node[label]at(1.0,-1.23){$R_m$};
\node[label]at(0,-2.6){表头：满偏电流 $I_g$，内阻 $R_g$};\end{scope}
\node[card,text width=6.5cm,draw=blue!65,fill=blue!4]at(3.9,1.1){$R_s=I_gR_g/(I-I_g)$};
\node[card,text width=6.5cm,draw=green!65,fill=green!4]at(3.9,-1.1){$R_m=U/I_g-R_g$};

\end{tikzpicture}\end{document}